\documentclass[12pt,a4paper]{article}

\usepackage{fontspec}
\setmainfont{CMU Serif}
\setsansfont{CMU Sans Serif}
\setmonofont{CMU Typewriter Text}
\usepackage[russian]{babel}

\usepackage{amsmath}
\usepackage{amssymb}
\usepackage{amsfonts}

\usepackage{graphicx}
\usepackage{geometry}
\geometry{top=3cm, bottom=2.5cm, left=3cm, right=2cm, headheight=1.5cm, headsep=0.8cm}

\usepackage{fancyhdr}
\pagestyle{fancy}
\fancyhf{}

\fancyhead[L]{\includegraphics[height=1.2cm]{../img/logo.png}}
\fancyhead[R]{\small\scshape STREMO: Математические модели}
\fancyfoot[C]{\thepage}
\renewcommand{\headrulewidth}{0.4pt}
\renewcommand{\footrulewidth}{0pt}

\usepackage{titlesec}
\titleformat{\section}{\Large\bfseries}{\thesection.}{0.5em}{}
\titleformat{\subsection}{\large\bfseries}{\thesubsection.}{0.5em}{}

\begin{document}

\thispagestyle{fancy}

\begin{center}
    \vspace{0.4cm}
    \Large{\textbf{Спецификация математических моделей и алгоритмов обработки}}\\
        \vspace{0.5cm}
    \rule{\textwidth}{0.4pt}
\end{center}

\vspace{1cm}

\section{1. Оптимизация размера медиа-чанка (Streaming Pipeline)}

Для балансировки задержки `Latency` и пропускной способности `Bandwidth` при передаче живого видеопотока по протоколу gRPC необходимо аналитически определить оптимальный размер медиа-чанка $S_{chunk}$ в битах.

Введем обозначения:
\begin{itemize}
    \item $B_{network}$ --- доступная пропускная способность канала связи (бит/с).
    \item $R_{bitrate}$ --- битрейт видеопотока (бит/с), при условии $B_{network} > R_{bitrate}$.
    \item $T_{overhead}$ --- константная задержка на сериализацию Protobuf и отправку сетевого кадра (с).
\end{itemize}

Длительность контента, содержащегося в одном чанке, вычисляется как $$D_{chunk} = \frac{S_{chunk}}{R_{bitrate}}$$.\\
Время физической передачи данного чанка по сети составляет $$T_{tx} = \frac{S_{chunk}}{B_{network}} + T_{overhead}$$.

Фундаментальное условие непрерывности трансляции (исключение опустошения буфера) требует, чтобы время передачи чанка не превышало его медийную длительность:
\begin{equation}
T_{tx} \le D_{chunk} \implies \frac{S_{chunk}}{B_{network}} + T_{overhead} \le \frac{S_{chunk}}{R_{bitrate}}
\end{equation}

Решая данное неравенство относительно искомого размера чанка $S_{chunk}$, получаем:
\begin{equation}
S_{chunk} \left( \frac{1}{R_{bitrate}} - \frac{1}{B_{network}} \right) \ge T_{overhead}
\end{equation}
\begin{equation}
S_{chunk} \ge \frac{T_{overhead} \cdot R_{bitrate} \cdot B_{network}}{B_{network} - R_{bitrate}}
\end{equation}

Поскольку end-to-end задержка прямо пропорциональна размеру чанка ($L \approx D_{chunk} + T_{tx}$), для ее минимизации размер чанка фиксируется на строгой нижней границе:
\begin{equation}
S_{opt} = \frac{T_{overhead} \cdot R_{bitrate} \cdot B_{network}}{B_{network} - R_{bitrate}}
\end{equation}

\section{2. Динамическая модель буферизации (ABR)}

Размер видеобуфера $Q(t)$ на стороне клиента в момент времени $t$ описывается обыкновенным дифференциальным уравнением:
\begin{equation}
\frac{dQ(t)}{dt} = \frac{C(t)}{R_v(t)} - 1
\end{equation}
где $C(t)$ --- текущая пропускная способность сети, а $R_v(t)$ --- битрейт выбранного слоя качества. Целевая функция алгоритма адаптивного битрейта (ABR) сводится к максимизации качества при жестком ограничении вероятности опустошения буфера:
\begin{equation}
\max_{R_v} \sum_{k=1}^{K} U(R_v(k)) - \lambda \sum_{k=1}^{K-1} |R_v(k+1) - R_v(k)|
\end{equation}
при условии $Q(t) > 0$, где $U$ --- функция полезности, а $\lambda$ --- весовой коэффициент штрафа за переключение качества.

\section{3. Матричные преобразования цветового пространства}

В процессе транскодирования видеопотока (переход YUV $\leftrightarrow$ RGB) применяется линейное преобразование цветовых метрик. В рамках математического аппарата системы STREMO вектор состояния пикселя $\mathbf{v}$ строго определяется как \textbf{вектор-столбец}:
\begin{equation}
\mathbf{v} = \begin{pmatrix} v_1 \\ v_2 \\ v_3 \end{pmatrix}
\end{equation}

Прямое преобразование из исходного цветового пространства задается операцией левостороннего умножения на матрицу преобразования $A$:
\begin{equation}
\mathbf{c} = A \mathbf{v}
\end{equation}

Для обратного вычисления исходных компонент применяется левостороннее умножение вектора-столбца $\mathbf{c}$ на обратную матрицу $A^{-1}$:
\begin{equation}
\mathbf{v} = A^{-1} \mathbf{c}
\end{equation}

Обратите внимание: операция транспонирования обратной матрицы в данном случае \textbf{не применяется}. Транспонирование необходимо исключительно при использовании конвенции векторов-строк (где $\mathbf{v}^T = \mathbf{c}^T (A^{-1})^T$). Так как система оперирует исключительно векторами-столбцами, форма $\mathbf{v} = A^{-1} \mathbf{c}$ сохраняет полную математическую корректность.

\end{document}